\exercisesection{线性代数的补充}

\begin{enumerate}
\def\labelenumi{\arabic{enumi})}
\tightlist
\item
  设
\end{enumerate}

\[
\mathbf {M} \rightarrow \mathbf {N} \xrightarrow {t} \mathbf {P} \rightarrow \mathbf {Q}
\]

\[
\begin{array}{c c c c c} \downarrow u & & \downarrow v & & \downarrow w \\ \downarrow & & \downarrow & & \downarrow s \end{array}
\]

\[
\mathrm{M} ^ {\prime} \rightarrow \mathrm{N} ^ {\prime} \xrightarrow {t ^ {\prime}} \mathrm{P} ^ {\prime} \rightarrow \mathrm{Q} ^ {\prime}
\]

A-模的一个交换图，其中各行都是正合的；假设 \(u\) 是满射且 \(s\) 是单射。证明：

\[
\operatorname{Ker} (w) = t (\operatorname{Ker} (v)) \quad \operatorname{Im} (v) = t ^ {\prime - 1} (\operatorname{Im} (w)).
\]

\begin{enumerate}
\def\labelenumi{\arabic{enumi})}
\setcounter{enumi}{1}
\tightlist
\item
  设 \(u: \mathbf{M} \to \mathbf{N}, v: \mathbf{N} \to \mathbf{P}\) 为 A-模的两个同态。证明存在正合列：
\end{enumerate}

\[
0 \rightarrow \operatorname{Ker} (u) \rightarrow \operatorname{Ker} (v \circ u) \rightarrow \operatorname{Ker} (v) \rightarrow \operatorname{Coker} (u) \rightarrow \operatorname{Coker} (v \circ u) \rightarrow \operatorname{Coker} v \rightarrow 0.
\]

\begin{enumerate}
\def\labelenumi{\arabic{enumi})}
\setcounter{enumi}{2}
\tightlist
\item
  在 A-模的图中：
\end{enumerate}

\[
\begin{array}{c} \mathrm{P} \xleftarrow {} \mathrm{M} \\ \Big \downarrow \quad \mathrm{R} ^ {\swarrow} \Big \downarrow \\ \mathrm{N} \xleftarrow {} \mathrm{Q} \end{array}
\]

图中提取出的两个三角形的交换性是否必然蕴含整个图的交换性？

\begin{enumerate}
\def\labelenumi{\arabic{enumi})}
\setcounter{enumi}{3}
\item
  设 \(k\) 为一个域， \(B = k[T]\) ，设 A 为子环 \(k[T^2, T^3] \subset k[T]\) 。证明 A-模 \(B\) 是无挠的，但不是平坦 A-模。
\item
  设 \(k\) 为一个域， \(\mathbf{C} = k[\mathbf{X}, \mathbf{Y}]\) ；考虑子环 \(\mathbf{A} = k[\mathbf{X}^4, \mathbf{Y}^4]\) , \(\mathbf{B} = k[\mathbf{X}^4, \mathbf{X}^3\mathbf{Y}, \mathbf{XY}^3, \mathbf{Y}^4]\) 与 \(\tilde{\mathbf{B}} = k[\mathbf{X}^4, \mathbf{X}^3\mathbf{Y}, \mathbf{X}^2\mathbf{Y}^2, \mathbf{XY}^3, \mathbf{Y}^4]\) ，它们都是 \(\mathbf{C}\) 的子环。包含关系 \(\mathbf{A} \subset \mathbf{B} \subset \tilde{\mathbf{B}}\) 使我们可以在 \(\mathbf{B}\) 与 \(\tilde{\mathbf{B}}\) 上定义 A-模结构。证明 \(\tilde{\mathbf{B}}\) 是平坦 A-模，而 A-模 \(\mathbf{B}\) 不是平坦的。
\item
  证明：一个 A-模 E 是平坦的，当且仅当它满足以下条件（参见 X，第 15 页，注记）：
\end{enumerate}

(ii'') 对于 A 中元素的每个有限族 \((c_{i})_{i\in I}\) ，方程的每个解 \(e = (e_i)_{i\in I}\) （方程 \(\sum_{i\in I}c_i e_i = 0\) 的解）都可以写成 \(b_1 z_1 + \cdots + b_n z_n\) ，其中 \(b_1, \ldots, b_n \in E\) ，并且对于 \(r = 1, \ldots, n\) , \(z_r = (z_{r,i})_{i \in I}\) 是方程 \(\sum_{i=1}^{n} z_{r,i} c_i = 0\) 的一个解。

\begin{enumerate}
\def\labelenumi{\arabic{enumi})}
\setcounter{enumi}{6}
\tightlist
\item
  设 E 是一个幺半群，S 是 E 的一个满足消去律且满足 III，第 121 页，习题 17 的假设的子半群，因而存在一个“分式幺半群” \(\overline{\mathbf{E}}\) 以及一个单同态 \(\varepsilon: \mathbf{E} \to \overline{\mathbf{E}}\) 。设 A 是一个交换环。由 \(\varepsilon\) 可导出幺半群代数的一个同态 \(u: \mathbf{A}^{(\mathbf{E})} \to \mathbf{A}^{(\overline{\mathbf{E}})}\) ，它使我们可以把 \(\mathbf{A}^{\overline{\mathbf{E}}}\) 视为左 \(\mathbf{A}^{(\mathbf{E})}\) -模。证明这个模在 \(\mathbf{A}^{(\mathbf{E})}\) 上是平坦的（将 \(\mathbf{A}^{\overline{\mathbf{E}}}\) 写成秩为 1 的自由 \(\mathbf{A}^{(\mathbf{E})}\) -模的滤过归纳极限）。
\end{enumerate}

* 8) 设 A 为区间 {[}0, 1{]} 上连续实值函数构成的环。

\begin{enumerate}
\def\labelenumi{\alph{enumi})}
\item
  证明在 0 处取值为零的函数构成的理想是平坦 A-模（证明它是理想 Af 的滤过归纳极限，其中 f 是仅在 0 处取值为零的函数）。
\item
  证明在 0 的某个邻域中取值为零的函数构成的理想 I 是投射 A-模（证明存在一个序列 \((f_n)_{n\geqslant 1}\) , \(f_n\in \mathbf{I}\) ，使得 \(\operatorname {Supp}(f_n)\subset \left[\frac{1}{n + 2},\frac{1}{n}\right]\) ，并且 \(\sum_{n}f_{n}(x) = 1\) 对 \(x\neq 0\) 成立；由此推出
\end{enumerate}

I 的每个元素 \(h\) 都可以写成 \(h = \sum_{n}^{t_{h}^{\prime}} (f_n h) e_n\) ，其中 \(e_n = \sum_{r \leqslant n+1}^{n} f_r\) ，并使用 II，第 46 页，命题 12。*

\begin{enumerate}
\def\labelenumi{\arabic{enumi})}
\setcounter{enumi}{8}
\tightlist
\item
  设 P 是一个 A-模。证明以下条件是等价的：
\end{enumerate}

α) 为使右 A-模的序列 \(M' \xrightarrow{u} M \xrightarrow{v} M''\) 正合，必要且充分条件是序列

\[
\mathbf {M} ^ {\prime} \otimes_ {\mathbf {A}} \mathbf {P} \xrightarrow {u \otimes 1} \mathbf {M} \otimes_ {\mathbf {A}} \mathbf {P} \xrightarrow {v \otimes 1} \mathbf {M} ^ {\prime \prime} \otimes_ {\mathbf {A}} \mathbf {P}
\]

是正合的。

\(\beta)\) P 是平坦的，并且对于每一个非零右 A-模 M，有 M \(\otimes_{A} P \neq (0)\) .

γ) P 是平坦的，且对所有右 A 模 M、N，典范同态

\[
\operatorname{Hom} _ {\mathbf {A}} (\mathbf {M}, \mathbf {N}) \rightarrow \operatorname{Hom} _ {\mathbf {Z}} (\mathbf {M} \otimes_ {\mathbf {A}} \mathbf {P}, \mathbf {N} \otimes_ {\mathbf {A}} \mathbf {P})
\]

是单射。

δ) P 是平坦的，并且对于 A 的每一个极大右理想 m，都有 mP ≠ P。

此时称 P 为忠实平坦 A-模。

\begin{enumerate}
\def\labelenumi{\arabic{enumi})}
\setcounter{enumi}{9}
\tightlist
\item
  \begin{enumerate}
  \def\labelenumii{\alph{enumii})}
  \tightlist
  \item
    证明一个忠实平坦模是忠实的（II，第28页）。
  \end{enumerate}
\end{enumerate}

\begin{enumerate}
\def\labelenumi{\alph{enumi})}
\setcounter{enumi}{1}
\item
  Z-模 Q 是忠实且平坦的，但不是忠实平坦的。
\item
  设 \((p_n)\) 为严格递增的素数序列，并设 A 为乘积环 \(\prod_{n} \mathbf{Z}/p_n \mathbf{Z}\) 。
\end{enumerate}

证明 \(Z/p_{n}Z\) 的直和 E 是忠实的投射 A-模，但不是忠实平坦的（注意 E 是 A 的一个理想，满足 \(E^{2}=E\) ）。

\begin{enumerate}
\def\labelenumi{\arabic{enumi})}
\setcounter{enumi}{10}
\item
  设 \(0 \to \mathbf{P}' \to \mathbf{P} \to \mathbf{P}'' \to 0\) 为 A-模的正合列。假设 \(\mathbf{P}'\) 与 \(\mathbf{P}''\) 是平坦的，且其中一个是忠实平坦的。证明 \(\mathbf{P}\) 是忠实平坦的。
\item
  设 \(\rho: A \to B\) 为一个环同态，使得 \(B\) 成为 \(A\) 上的忠实平坦右模；设 \(M\) 为一个 \(A\) 模。
\end{enumerate}

\begin{enumerate}
\def\labelenumi{\alph{enumi})}
\item
  证明：M 是有限生成的（相应地，有限表示的），当且仅当 B-模 \(\mathbf{B} \otimes_{\mathbf{A}} \mathbf{M}\) 亦然。
\item
  我们假设 \(\rho(A)\) 包含在 B 的中心中。证明：M 是平坦的（相应地，忠实平坦的，相应地，有限生成投射的），当且仅当 B-模 \(B \otimes_{A} M\) 是如此。
\end{enumerate}

\begin{enumerate}
\def\labelenumi{\arabic{enumi})}
\setcounter{enumi}{12}
\tightlist
\item
  证明：整环 A 上每个有限生成的平坦模 P 都是投射的（考虑一个表示 \(0 \rightarrow R \xrightarrow{L} L \xrightarrow{P} 0\) ，其中 L 是有限生成的自由模；选取 R 的一个有限生成子模 \(R'\) ，使得 \(R/R'\) 是挠模，并应用第 14 页定理 1）。
\end{enumerate}

¶ 14) a) 设 \(0 \to \mathbf{R} \to \mathbf{L} \to \mathbf{E} \to 0\) 为 A-模的一个正合列，其中 L 是自由 A-模；设 \((e_{a})\) 为 L 的一组基。证明以下条件是等价的：

β) 对每个 \(x \in \mathbb{R}\) ，若 \(\mathfrak{a}_{x}\) 是 \(x\) 在基 \((e_{\alpha})\) 上的分量生成的右理想，则有 \(x \in \mathfrak{a}_{x} \mathbb{R}\) 。

\(\gamma)\) 对每一个 \(x \in \mathbb{R}\) ，存在同态 \(u_x: \mathbb{L} \to \mathbb{R}\) 使得 \(u_x(x) = x\) 。

δ) 对 R 的元素的每个有限序列 \((x_{i})_{1\leqslant i\leqslant n}\) ，存在同态 \(u:L\to R\) 使得 \(u(x_{i})=x_{i}\) 对 \(1\leqslant i\leqslant n\) 成立（利用第 14 页定理 1。）

\begin{enumerate}
\def\labelenumi{\alph{enumi})}
\setcounter{enumi}{1}
\item
  设 r 为 A 的根，并设 \(0 \rightarrow R \rightarrow L \rightarrow E \rightarrow 0\) 为 A-模的一个正合列，其中 L 是自由的。假设 E 是平坦的，且 R 包含于 rL 中。证明此时 R = 0（用 a) 的记号，注意 \(a_{x}\) 是有限生成的理想，并且我们有 \(a_{x} = a_{x} r\) ）。
\item
  设 E 是一个平坦的有限生成 A-模；假设存在 A 的一个双边理想 b 包含于 A 的根中，使得 E/bE 是自由 (A/b)-模。证明此时 E 是自由 A-模（注意到存在有限生成自由 A-模 L 使得 L/bL 同构于 E/bE，并应用 b））。特别地，局部环上的每个有限生成平坦模都是自由的。
\end{enumerate}

\begin{enumerate}
\def\labelenumi{\arabic{enumi})}
\setcounter{enumi}{14}
\tightlist
\item
  设 A 为一个环， \(a\) 为 A 的一个元素。证明以下性质是等价的： \(\alpha)\) \(a\in aAa\) 。
\end{enumerate}

β) Aa 是模 \(A_{s}\) 的一个直和项。

γ) \(A_{s}/Aa\) 是平坦的 A-模。

δ) 对于 A 的每个右理想 b，我们有 b ∩ Aa = ba。

\begin{enumerate}
\def\labelenumi{\arabic{enumi})}
\setcounter{enumi}{15}
\tightlist
\item
  设 A 是一个环。证明以下性质是等价的：
\end{enumerate}

α) 每个元素 \(a \in A\) 满足习题15的等价性质。

β) 每个有限生成的左理想都是 \(A_{s}\) 的直和项。

γ) 每个左 A-模都是平坦的。

δ) 每个右A-模都是平坦的。

我们于是称 A 为绝对平坦环（参见 VIII，§ 8，习题 15）。

\begin{enumerate}
\def\labelenumi{\arabic{enumi})}
\setcounter{enumi}{16}
\tightlist
\item
  设 A 是一个绝对平坦环（习题 16）。
\end{enumerate}

\begin{enumerate}
\def\labelenumi{\alph{enumi})}
\item
  设 P 为投射 A-模。证明 P 的每个有限生成子模都是 P 的直和项（化归到 P 为有限生成自由模的情形）。
\item
  证明每一个投射A-模都是循环子模的直和，且这些循环子模同构于A的循环理想。（利用Kaplansky定理（II，第183页，习题2）将情形化归为P具有可数生成族的情形，然后利用a）。）
\item
  证明 A 的每个由可数生成元族生成的理想都是投射的。
\item
  给出一个绝对平坦环 A 以及一个不是投射模的有限生成 A-模的例子。
\end{enumerate}

18) 设 M 是一个 A-模。证明以下条件是等价的：

α) M 是有限生成的（相应地，有限表示的）。

β) 对于相对于滤过预序集 I 的每个 A-模归纳系 \((\mathbf{N}_{i}, u_{ji})\) ，典范同态 \(\varinjlim\operatorname{Hom}_{\mathbf{A}}(\mathbf{M}, \mathbf{N}_{i}) \to \operatorname{Hom}_{\mathbf{A}}(\mathbf{M}, \varinjlim \mathbf{N}_{i})\) 是单射（相应地，双射）。

γ) 对于右 A-模的每一个族 \((\mathbf{P}_{j})_{j\in J}\) ，典范同态

\[
\left(\prod_ {j \in J} P _ {j}\right) \otimes_ {A} M \rightarrow \prod_ {j \in J} \left(P _ {j} \otimes_ {A} M\right)
\]

（II，第61页）是满射（相应地，双射）。

δ) 对于每个集合 I，典范同态 \(A_{d}^{1} \otimes_{A} M \to M^{1}\) 是满射（相应地，双射）。 19) 设 A 和 B 是两个环， \(\varphi : A \to B\) 为一个同态，M 是一个 B-模。我们赋予 \(B \otimes_{A} M\) 与 \(\operatorname{Hom}_{A}(B, M)\) 由 B 的 B-模结构诱导的左 B-模结构。称 M 相对于 A 是投射的，或 \(\varphi\) -投射的，如果典范满射 \(B \otimes_{A} M \to M\) 具有 B-线性截面；称 M 相对于 A 是内射的，或 \(\varphi\) -内射的，如果典范单射

\[
\mathbf {M} \rightarrow \operatorname{Hom} _ {\mathbf {A}} (\mathbf {B}, \mathbf {M})
\]

有一个 B-线性收缩。

\begin{enumerate}
\def\labelenumi{\alph{enumi})}
\item
  如果 M 作为 A-模是投射的（相应地，内射的）且是 \(\varphi\) -投射的（相应地， \(\varphi\) -内射的），则 M 是投射的（相应地，内射的）B-模。
\item
  对于每个 A-模 N，B-模 \(B \otimes_{A} N\) （相应地， \(\operatorname{Hom}_{A}(B, N)\) ）是 \(\varphi\) -投射的（相应地， \(\varphi\) -内射的）。
\item
  我们假设 \(\varphi(A)\) 包含在 B 的中心中。设 N 是一个 A-模，P 是一个 \(\varphi\) -投射的 B-模。证明 B-模 \(P \otimes_{A} N\) （相应地，右 B-模 \(\operatorname{Hom}_{A}(P, N)\) ）是 \(\varphi\) -投射的（相应地， \(\varphi\) -内射的）。
\end{enumerate}

\begin{enumerate}
\def\labelenumi{\arabic{enumi})}
\setcounter{enumi}{19}
\item
  设 A 是一个主理想环， \(a\) 为 A 的一个非零元素，B = A/Aa。证明 B-模 B \(_{s}\) 是内射的。
\item
  证明以下条件是等价的：
\end{enumerate}

α) A 是诺特环。

β) 每个作为内射 A-模的归纳极限的 A-模都是内射的。

γ) 每个作为内射 A-模的直和的 A-模都是内射的。

（为了证明 \(\gamma\) ）蕴含 \(\alpha\) ），设 \((\mathfrak{a}_n)_{n\geqslant 1}\) 为 A 的左理想组成的一个递增序列， \(\mathfrak{a} = \bigcup_{n}\mathfrak{a}_{n}\) ,

\(I_{n}\) 为 \(\mathfrak{a}/\mathfrak{a}_{n}\) 的一个内射包络。从 \(\mathfrak{a}\) 到 \(\oplus I_{n}\) 的自然同态是某个同态 \(f: A \to \oplus I_{n}\) 的限制；考虑 \(f(1)\) 。

\begin{enumerate}
\def\labelenumi{\arabic{enumi})}
\setcounter{enumi}{21}
\tightlist
\item
  设 A 为主理想环，K 为其分式域， \(\pi\) 为 A 的一个极值元素。证明 A-模 K/A 的 \(\pi\) -准素分量是 A-模 A/ \(\pi^n\) A 的一个内射包络，对每一个 \(n \geqslant 1\) 。 ¶ 23) 设 I 是一个内射 A-模，E 为环 End \(_A\) (I)，r 为 E 中使得 I 是 Ker( \(u\) ) 的内射包络的元素 \(u\) 组成的集合。
\end{enumerate}

\begin{enumerate}
\def\labelenumi{\alph{enumi})}
\item
  证明 r 是一个双边理想，并且环 E/r 是绝对平坦的（习题 16）。（设 u 是 E 的一个元素，N 是 I 的一个子模，使得 N ∩ Ker u = 0 且对此性质是极大的；证明存在 v ∈ E 使得对所有 x ∈ N 有 vu(x) = x，并由此推出 uvu - u ∈ r。）
\item
  证明 r 是 E 的根（包含关系 r(E) ⊂ r 由 a) 得出）；证明对于每个 u ∈ r，1 - u 是双射的）。
\item
  假设 I 是有限个不可分解内射模的直和。证明环 E/r 是半单的。
\end{enumerate}

设 \(i: \mathbf{A} \to \mathbf{I}\) 为 A-模 \(\mathbf{A}_{s}\) 的一个内射包络。我们设 \(\mathbf{E} = \operatorname{End}_{\mathbf{A}}(\mathbf{I})\) 与 \(\mathbf{B} = \operatorname{End}_{\mathbf{E}}(\mathbf{I})\) ，使得 \(\mathbf{B}\) 是 A-模 I 的双交换子（VIII，§ 3，定义 1）。我们把 \(\mathbf{A}\) 视为 \(\mathbf{B}\) 的一个子环。用 \(j: \mathbf{B} \to \mathbf{I}\) 与 \(p: \mathbf{E} \to \mathbf{I}\) 表示由 \(j(b) = b(i(1))\) 与 \(p(u) = u(i(1))\) （对 \(b \in \mathbf{B}\) , \(u \in \mathbf{E}\) ）定义的映射。

\begin{enumerate}
\def\labelenumi{\alph{enumi})}
\item
  证明 \(p\) 是满射且 \(j\) 是单射。
\item
  证明以下条件是等价的：
\end{enumerate}

α) 映射 p 是双射的。

β) 映射 j 是双射。

γ) A-模 B 是内射的。

若这些条件得到满足，则映射 \(j^{-1} \circ p\) 是从环 E 到 B 的相反环上的一个同构。

\begin{enumerate}
\def\labelenumi{\alph{enumi})}
\setcounter{enumi}{2}
\item
  设 \(\mathfrak{F}\) 为左理想 \(a\) （ \(A\) 中）组成的集合，满足 \(a \cap b \neq (0)\) 对每个非零左理想 \(b\) 都成立。用 \(s(A)\) 表示满足如下条件的元素 \(a\) （ \(A\) 中）组成的集合：存在 \(a \in \mathfrak{F}\) 使得 \(aa = 0\) 。证明 \(s(A)\) 是 \(A\) 的一个双边理想。
\item
  证明以下条件是等价的：
\end{enumerate}

α) \(s(\mathbf{A}) = 0.\)

β) 环 E 是半本原的。

γ) 环 B 是绝对平坦的。

若这些条件满足，则 \(b\) ) 的等价性质成立。

（设 \(s(I)\) 为元素 \(x\) （I 中）组成的集合，其中存在 \(\mathfrak{a} \in \mathfrak{F}\) 使得 \(\mathfrak{a} x = 0\) 。利用习题 23，证明 \(p^{-1}(s(I))\) 等于 E 的根 \(r\) ，由此得到蕴含关系 \(\beta) \Rightarrow \alpha\) 。然后证明 \(\alpha\) 蕴含 \(r = \text{Ker } (p) = 0\) ，它蕴含 \(\beta\) 与 \(\gamma\) （根据习题 23），也蕴含条件 \(\alpha\) （ \(b\) ) 的）。最后证明 \(\gamma\) 蕴含 \(\alpha\) 。）

\begin{enumerate}
\def\labelenumi{\alph{enumi})}
\setcounter{enumi}{4}
\item
  若 A 是左诺特的，则理想 s(A) 是幂零的（首先证明 s(A) 的每个元素 a 都是幂零的，考虑 \(a^{n}\) 的左零化子；然后利用 VIII，§ 9，习题 8 b））。¶ 25) 假设环 A 是左诺特的且不含非零幂零双边理想。a) 证明习题 24 中定义的环 B 是半单的（“Goldie 定理”；利用习题 23 c）以及习题 24 d）和 e））。若此外 (0) 是 A 中的素双边理想（VIII，§ 8，习题 19），证明 B 是单的。
\item
  将 A 等同于 B 的一个子环。证明以下性质：
\end{enumerate}

\begin{enumerate}
\def\labelenumi{(\roman{enumi})}
\item
  A 中每个左可消去的元素在 B 中都可逆（从而在 A 中可消去）。
\item
  每个元素 \(b\) （ \(\mathbf{B}\) 中）都可以写成 \(s^{-1}a\) ，其中 \(s \in \mathbf{A}\) , \(a \in \mathbf{A}\) ，且 \(s\) 可消去。
\end{enumerate}

有时称 B 为 A 的左全分式环。

（为证明 (ii)，归结到 B 是单的情形。然后证明 A 的理想 (0) 是素的。设 a 是 F 的一个理想，a 是 a 中使得右理想 aB 在理想 a'B（a' ∈ a）中极大的元素，x 是 B 中满足 axa = a 的元素。证明 B(1 - xa) ⊂ aB，由此

\[
(\mathbf {B} (1 - a x) \cap \mathfrak {a}) \cdot (\mathbf {B} (1 - x a) \cap \mathfrak {a}) = (0);
\]

推出 \(a\) 可消去。最后取 a 为所有 \(a \in A\) 使得 \(ab \in A\) 的元素组成的集合来得出结论。）c) 假设 \(A\) 的每个非零元素都是可消去的。证明 \(B\) 是一个域，它是 \(A\) 的左分式域（参见第一卷，第 155 页，习题 15）。

¶ 26) a) 证明以下条件是等价的：

α) 环 A 是右诺特的，且 A-模 \(A_{s}\) 是内射的。

\(\alpha^{0}\) A 是左诺特的且 A-模 \(A_{d}\) 是内射的。

β) 环 A 是自内射的（VIII，§ 2，习题 11）。

（证明 \(\alpha\) ) 蕴含每个右理想 r 满足 \(r = r^0\) 。然后证明 \(r(A)\) 是幂零的，并利用习题 23 推出 A 是右阿廷环。接着证明单左 A-模的对偶是单模。由此推出 A 是左诺特环，并且最终是自内射的。b) 证明环 A 是自内射的，当且仅当每个内射 A-模都是投射的且每个投射 A-模都是内射的。

¶ 27) 假设环 A 是左诺特的。

\begin{enumerate}
\def\labelenumi{\alph{enumi})}
\item
  设 I 是一个不可分解的内射 A-模。证明 I 的非零子模的零化子理想组成的集合具有一个最大元 \(\mathfrak{A}(I)\) ，它是一个素双边理想（VIII，§ 8，习题 19）。
\item
  证明：对于 A 的每一个素双边理想 p，存在一个不可分解的内射 A-模 I，使得 \(\mathfrak{A}(I) = \mathfrak{p}\) （考虑 A/p 的内射包络的一个不可分解直因子）。
\item
  假设 A 是交换的。证明每个不可分解的内射 A-模 I 都是 A/U(I) 的内射包络。由此推出映射 I ↦ U(I) 定义了从不可分解内射 A-模的同构类集合 S 到 A 的素理想集合上的一个双射。逆双射把理想 p 对应于 A/p 的内射包络。
\item
  证明：为使映射 I ↦ \(\mathfrak{A}(I)\) （从 \(\mathfrak{A}\) 到 A 的素双边理想集合上）是双射，只需 A 满足以下条件 (H)：
\end{enumerate}

\begin{enumerate}
\def\labelenumi{(\Alph{enumi})}
\setcounter{enumi}{7}
\tightlist
\item
  对于 A 的每一个左理想 a，存在 A/a 的元素 \(x_{1}, \ldots, x_{k}\) ，使得
\end{enumerate}

\[
\operatorname{Ann} (\mathbf {A} / \mathfrak {a}) = \bigcap_ {i} \operatorname{Ann} (x _ {i}).
\]

\begin{enumerate}
\def\labelenumi{\alph{enumi})}
\setcounter{enumi}{4}
\tightlist
\item
  证明在以下情形中条件 (H) 得到满足：
\end{enumerate}

α) A 的所有左理想都是双边理想。

β) A 是阿廷的。

γ) A 的中心 Z 是诺特环，且 A 是有限生成的 Z-模。

\begin{enumerate}
\def\labelenumi{\alph{enumi})}
\setcounter{enumi}{5}
\tightlist
\item
  设 k 为一个域， \(Q^{+}\) 为满足 \(\geqslant 0\) 的有理数组成的加法幺半群，A 为代数 \(k^{(\mathbf{Q}^{+})}\) ，a 为 A 的一个非零元素，I 为 A/Aa 的内射包络。证明 A-模 I 是不可分解的，但不包含与 A/p 同构的子模，其中 p 为素理想（注意 A 的理想集是全序的，且 A 仅包含两个素理想）。
\end{enumerate}

\begin{enumerate}
\def\labelenumi{\arabic{enumi})}
\setcounter{enumi}{27}
\tightlist
\item
  设 A 为诺特环，M 为 A-模，I 为 M 的内射包络，它是某个族 \((I_{\alpha})_{\alpha \in J}\) 的不可分解子模的直和。我们用 Ass(M) 表示素双边理想 \(\mathfrak{U}(I_{\alpha})\) （ \(\alpha \in J\) ）组成的集合（习题 27）；若 \(p \in \operatorname{Ass}(M)\) ，则称 p 与 M 相伴。
\end{enumerate}

\begin{enumerate}
\def\labelenumi{\alph{enumi})}
\item
  若 M ≠ 0，则 Ass (M) ≠ ∅；若 M 是有限生成的，则集合 Ass (M) 是有限的。
\item
  证明：素双边理想 p 与 M 相伴当且仅当 M 包含一个子模 N，使得对 N 的每个非零子模 N'，都有 Ann (N') = p。若 A 是交换的，则 p 与 M 相伴当且仅当 M 包含一个与 A/p 同构的子模。
\item
  设 M' 是 M 的一个子模。证明我们有
\end{enumerate}

\[
\operatorname{Ass} \left(\mathbf {M} ^ {\prime}\right) \subset \operatorname{Ass} (\mathbf {M}) \subset \operatorname{Ass} \left(\mathbf {M} ^ {\prime}\right) \cup \operatorname{Ass} \left(\mathbf {M} / \mathbf {M} ^ {\prime}\right).
\]

\begin{enumerate}
\def\labelenumi{\alph{enumi})}
\setcounter{enumi}{3}
\tightlist
\item
  我们设 \(J_{p} = \sum_{\mathfrak{M}(I_{\alpha}) = p} I_{\alpha}\) 与 \(Q(p) = \left( \sum_{q \neq p} J_{q} \right) \cap M\) 。证明我们有 \(\bigcap_{p \in \operatorname{Ass}(M)} Q(p) = 0\) ，并且对 Ass(M) 的每个异于 Ass(M) 的子集 E，有 \(\bigcap_{p \in E} Q(p) \neq 0\) 。对每个 \(p \in \operatorname{Ass}(M)\) ，我们有
\end{enumerate}

Ass \((\mathbf{M} / \mathbf{Q}(\mathfrak{p})) = \{\mathfrak{p}\}\) .

\begin{enumerate}
\def\labelenumi{\alph{enumi})}
\setcounter{enumi}{4}
\tightlist
\item
  设 a 为 A 的中心的一个元素。证明：以 a 为比率的位似变换在 M 中是单射的，其必要且充分条件是 a 不属于与 M 相伴的任何一个素理想。
\end{enumerate}

假设环 A 是交换的且是诺特的。设 a 为 A 的一个理想。

\begin{enumerate}
\def\labelenumi{\alph{enumi})}
\tightlist
\item
  设 H 为一个 A-模，使得 H 的每个元素都被 a 的某个幂零化。证明：要使 H 是内射的，只需对每个整数 \(n \geqslant 1\) 、每个被 \(a^{n}\) 零化的有限生成 A-模 M，以及 M 的每个子模 \(M'\) ，自然同态 \(\operatorname{Hom}_{\mathbf{A}}(\mathbf{M}, \mathbf{H}) \to \operatorname{Hom}_{\mathbf{A}}(\mathbf{M}', \mathbf{H})\) 是满射。
\end{enumerate}

（利用如下事实：对 A 的每个理想 b，存在一个整数 \(k \geqslant 0\) 使得 \(\mathfrak{b} \cap \mathfrak{a}^k \subset \mathfrak{ba}\) （AC，III，§ 3，第 1 号，推论 2）。）

\begin{enumerate}
\def\labelenumi{\alph{enumi})}
\setcounter{enumi}{1}
\item
  设 M 为一个内射 A-模。对于 \(n \geqslant 1\) ，我们用 \(\mathbf{H}_{n}\) 表示 M 中由被 \(\alpha^{n}\) 零化的元素组成的子模；令 \(\mathbf{H} = \bigcup \mathbf{H}_{n}\) 。证明 H 是一个内射 A-模。
\item
  证明 A/a 的内射包络的每个元素都被 a 的某个幂零化。* 30) 假设环 A 是局部、交换且诺特的；用 m 表示其极大理想，并令 k = A/m。设 I 为一个 A-模，使得 I 的每个元素都被 m 的某个幂零化。对每个 A-模 M，用 T(M) 表示 A-模 Hom\_A (M, I)，并用 c\_M : M → T(T(M)) 表示由
\end{enumerate}

\[
\left(c _ {\mathbf {M}} (m)\right) (u) = u (m) \quad \text { for } \quad m \in \mathbf {M}  , \quad u \in \mathrm{T} (\mathbf {M})  .
\]

定义的同态。证明下列性质是等价的：

α) 对每个有限长度 A-模 M，A-模 T(M) 是有限长度的，且同态 \(c_{M}: M \to T(T(M))\) 是双射。

β) 对每个有限长度的 A-模 M，有 long T(M) = long M。

γ) A-模 I 是内射的，且 A-模 k 与 T(k) 同构。

δ) A-模 I 与 k 的内射包络同构。

当这些条件满足时，有时称 I 为环 A 的一个对偶化模。（证明性质 γ) 与其他三个性质中的每一个等价：为看出 α) 或 β) 蕴含 γ)，使用习题 29。）

\begin{enumerate}
\def\labelenumi{\arabic{enumi})}
\setcounter{enumi}{30}
\item
  在上一习题的假设下，设 B 为一个局部、交换、诺特环，并设 \(\rho: A \to B\) 为一个环同态，使得 B 成为有限生成 A-模。若 I 是 A 的一个对偶化模，证明 B-模 \(\mathrm{Hom}_{\mathbf{A}}(\mathbf{B}, I)\) 是 B 的对偶化模。特别地，对 A 的每个理想 \(a\) ，I 中由被 \(a\) 零化的元素组成的子模是 A/a 的一个对偶化模。
\item
  \begin{enumerate}
  \def\labelenumii{\alph{enumii})}
  \tightlist
  \item
    我们保留习题 30 的假设；进一步假设 A 包含一个域 \(k_0\) 使得 \(k\) 是有限维的 \(k_0\) -向量空间。证明 A-模 \(l = \varinjlim_{n} \operatorname{Hom}_{k_0}(A/m^n, k_0)\)
  \end{enumerate}
\end{enumerate}

是 A 的对偶化模。

\begin{enumerate}
\def\labelenumi{\alph{enumi})}
\setcounter{enumi}{1}
\tightlist
\item
  设 \(k\) 为特征为零的域，B 为形式幂级数环 \(k[[\mathrm{T}_{1}, \ldots, \mathrm{T}_{n}]]\) 。我们赋予群 \(\mathbf{I} = k[\mathrm{X}_{1}, \ldots, \mathrm{X}_{n}]\) 一个 B-模结构，通过设定 \(\mathrm{T}_{i} \cdot \mathbf{P} = \partial \mathbf{P} / \partial \mathbf{X}_{i}\) 对所有 \(\mathbf{P} \in k[\mathrm{X}_{1}, \ldots, \mathrm{X}_{n}]\) 。证明 I 是 \(k\) 的一个内射包络。
\end{enumerate}

